\chapter{绪论}
\label{chap:introduction}

\section{研究背景}

在这里介绍课题背景、研究价值和实际需求。首次出现的缩略语应给出中文名称、英文名称和缩写；后文保持同一术语与符号写法。需要补充而不宜放入正文的说明，可使用脚注。\footnote{这是脚注示例，使用带圈编号。学校规范建议采用脚注，但未统一规定编号样式。}

\section{研究现状}

按主题归纳已有工作，并在正文中使用 \verb|\cite{...}| 引用文献。例如，\cite{lamport1994latex,knuth1984literate} 分别给出了文档排版与可复现技术写作的经典讨论。不要把参考文献仅堆在段末，应说明它与当前论述的关系。

会议论文和网络资料的著录示例分别见文献\cite{vaswani2017attention,latexproject}。此处只演示不同文献类型的引用格式，实际写作时应替换为与研究内容相关的来源。

\section{本文工作与结构}

说明本文解决的问题、采用的技术路线和各章安排。可以用列表概括研究工作：

\begin{enumerate}[label=(\arabic*)]
  \item 明确研究对象、问题边界和评价指标；
  \item 说明方法、数据或实验设计；
  \item 给出结果、讨论局限性并形成结论。
\end{enumerate}

本模板在第~\ref{chap:method} 章给出公式、单图和定理的写法，在第~\ref{chap:results} 章给出三线表、表注、宽表、子图与算法的写法。删除不适用于本人的示例，不要在最终稿保留占位文字。
